\documentclass[10pt,a4paper]{amsart}
\usepackage[margin=19mm]{geometry}
\usepackage{amsmath,amssymb,mathtools}
\usepackage[scr=rsfs,cal=euler]{mathalfa}
\usepackage{xfrac}
\usepackage[unicode,hidelinks,bookmarks=false]{hyperref}
\newcommand{\ie}{i.e.\ }
\newcommand{\cf}{cf.\ }
\newcommand{\resp}{respectively\ }
\newcommand{\Subsection}{\subsection}
\providecommand{\nobreakdash}{\nobreak\hbox{-}}
\setlength{\emergencystretch}{2em}
\multlinegap0pt
\usepackage{amssymb}
\usepackage{mathtools}
\usepackage[shortlabels]{enumitem}
\newtheorem{theorem}{Theorem}
\title{Existence of Pure Strategy Nash Equilibria in Finite Noncooperative Games}
\author{Shravan Luckraz}
\date{Source: arXiv:2606.26564v3 (CC BY 4.0)}
\begin{document}
\maketitle

\noindent\emph{Source and licence.} Existence of Pure Strategy Nash Equilibria in Finite Noncooperative Games. Version arXiv:2606.26564v3 (CC BY 4.0), distributed by arXiv under the Creative Commons Attribution 4.0 International licence, http://creativecommons.org/licenses/by/4.0/, as recorded in the arXiv licence metadata for this exact version and shown on the abstract page https://arxiv.org/abs/2606.26564v3. Full licence notice: \texttt{licenses/arxiv-2606.26564-CC-BY-4.0.txt}.\par\smallskip

\noindent\textit{Editorial scope.} Counted unit: the article's theorem theorem (source0.tex lines 1131--1137). The article's complete original proof of it is delivered with it (source0.tex lines 1139--1228, one \texttt{proof} environment of 90 lines). No mathematics is written, repaired or re-derived: the statement and the proof are the article's own lines, in the article's own notation, and no retained line is substituted or re-flowed.  No further source-local context is needed: every cross-reference of every retained line resolves inside the two delivered spans.  Nothing was omitted from any retained span, so the package carries no omission record.  No retained line contains a citation, so the wrapper carries no bibliography and no bibliography entry.  No retained line contains a figure environment, an image-inclusion command or drawing code, so no image is delivered and the package declares no asset dependency.  Editorial formatting only, in addition: an AMS \texttt{amsart} preamble that restates the article's own declarations and macros used by the retained lines, namely the environments \texttt{theorem} on separate counters, together with the standard packages those lines need.  The retained environments print in this document as Theorem 1 in order of appearance, whereas the article prints the same items with the numbering of its own full document.  The complete native source of the article as distributed in the arXiv e-print for this exact version is bundled unchanged as \texttt{sources/203697/source0.tex}.
\begin{theorem}
Let \(\Gamma=(N,\{S_i\},\{u_i\})\) be a finite game with \(M=|S|\). There
exists \(\Gamma'=(N,\{S_i\},\{u_i'\})\) argmax payoff equivalent to \(\Gamma\)
such that \(U'(s)=\sum_i u_i'(s)\) is injective. Moreover \(\Gamma'\) can be
constructed by perturbations of magnitude at most \(\varepsilon\) for any
chosen \(\varepsilon>0\) sufficiently small relative to payoff gaps.
\end{theorem}

\begin{proof}
Because \(S\) is finite, for each player \(i\) and each opponents' profile
\(s_{-i}\) the finite set \(\{u_i(x,s_{-i}): x\in S_i\}\) has a finite set of
pairwise differences. Define the minimal positive gap for player \(i\) at
\(s_{-i}\) by


\[
\gamma_{i,s_{-i}} := \min\{ |u_i(x,s_{-i})-u_i(x',s_{-i})| : x,x'\in S_i,\ x\ne x'\}.
\]


If all payoffs coincide for some \(i,s_{-i}\) then set \(\gamma_{i,s_{-i}}=+\infty\). Let


\[
\gamma := \min_{i,s_{-i}:\ \gamma_{i,s_{-i}}<\infty} \gamma_{i,s_{-i}}.
\]


If every \(\gamma_{i,s_{-i}}=+\infty\) ( since all payoffs identical for each
player given opponents), then any arbitrarily small perturbation that
separates totals will preserve best replies. This is true because there are no strict best
reply distinctions to preserve. Otherwise \(\gamma>0\).

Choose \(\delta\) such that \(0<\delta<\gamma/4\). We will construct
perturbations \(\epsilon_i(s)\) with \(|\epsilon_i(s)|<\delta\) for all
\(i,s\), and then adjust one coordinate per profile by at most \(\delta\) to
achieve distinct totals.

\textbf{Stage 1: small perturbations preserving best replies.} For each
profile \(s\) and each player \(i\) set \(\epsilon_i(s)\) arbitrarily with
\(|\epsilon_i(s)|<\delta\). Define provisional payoffs
\(\tilde u_i(s)=u_i(s)+\epsilon_i(s)\). For any fixed opponents' profile
\(s_{-i}\) and any two actions \(x,x'\in S_i\),


\[
\tilde u_i(x,s_{-i})-\tilde u_i(x',s_{-i})
= u_i(x,s_{-i})-u_i(x',s_{-i}) + \epsilon_i(x,s_{-i})-\epsilon_i(x',s_{-i}).
\]


The perturbation difference \(|\epsilon_i(x,s_{-i})-\epsilon_i(x',s_{-i})|\)
is at most \(2\delta<\gamma/2\). If \(u_i(x,s_{-i})-u_i(x',s_{-i})\ge \gamma\)
then the sign of the difference is preserved. If the original difference is
zero, the perturbation may break ties; this is acceptable because we only
need to preserve argmax sets, breaking ties arbitrarily is allowed as long
as we do not invert strict inequalities. Because \(\delta<\gamma/4\), any
original strict inequality with margin at least \(\gamma\) remains strict in
\(\tilde u\). Thus for every \(i\) and \(s_{-i}\), the set of maximizers of
\(\tilde u_i(\cdot,s_{-i})\) equals the set of maximizers of \(u_i(\cdot,s_{-i})\).
If some original differences are smaller than \(\gamma\) but nonzero, we
can refine the choice of \(\delta\) using the actual minimal positive gap
over all players and opponent profiles; the finite nature of the game makes
this possible.

\textbf{Stage 2: make totals distinct.} Let \(\tilde U(s)=\sum_i \tilde u_i(s)\).
We want to adjust the provisional payoffs by small amounts to obtain final
payoffs \(u_i'(s)\) with distinct totals \(U'(s)\). Enumerate profiles
\(s^{(1)},\dots,s^{(M)}\). Choose target totals \(a_1<\dots<a_M\) such that
\(|a_k-\tilde U(s^{(k)})|<\delta\) for all \(k\) (this is possible because
we can choose \(a_k=\tilde U(s^{(k)})+\theta_k\) with small distinct
\(\theta_k\) satisfying \(|\theta_k|<\delta\) and \(\theta_k\) strictly
increasing in \(k\)). For each profile \(s^{(k)}\) adjust the payoff of a
single designated player (say player 1) by setting


\[
u_1'(s^{(k)}) = \tilde u_1(s^{(k)}) + (a_k-\tilde U(s^{(k)})),
\]


and for \(i\ne 1\) set \(u_i'(s^{(k)})=\tilde u_i(s^{(k)})\). Because
\(|a_k-\tilde U(s^{(k)})|<\delta\), the adjustment to player 1's payoff is
less than \(\delta\) in magnitude. Therefore for any fixed opponents'
profile \(s_{-1}\), the ordering of \(u_1'(\cdot,s_{-1})\) remains the same
as that of \(\tilde u_1(\cdot,s_{-1})\) because the per‑action adjustments
are smaller than the minimal gap used in Stage 1. Hence argmax sets are
preserved for player 1 as well.

Thus the final payoffs \(u_i'\) satisfy:
First, for every \(i\) and \(s_{-i}\), \(\arg\max_{x\in S_i} u_i'(x,s_{-i})
    = \arg\max_{x\in S_i} u_i(x,s_{-i})\). (Argmax payoff equivalence.) Second, the totals \(U'(s)=\sum_i u_i'(s)\) equal the distinct targets \(a_k\),
    hence \(U'\) is injective. Third, all perturbations are bounded by \(2\delta<\gamma/2\), which can be
    made arbitrarily small by choosing \(\delta\) small relative to the
    minimal positive payoff gap \(\gamma\).

This completes the constructive proof.
\end{proof}

\end{document}
